\documentclass[12pt]{article}
\usepackage[utf8]{inputenc}
\usepackage{amsmath,amssymb,hyperref,array,xcolor,multicol,verbatim,mathpazo}
\usepackage[normalem]{ulem}
\usepackage[pdftex]{graphicx}
\usepackage{fullpage}
\usepackage{import}
\usepackage{adjustbox}
\usepackage{booktabs}
\usepackage{bbm}
\usepackage[font=normalsize,labelfont=bf]{caption}
%\captionsetup{justification=raggedright,singlelinecheck=false}
\usepackage{subcaption}
\usepackage{pdflscape}
\usepackage{threeparttable}
\usepackage{makecell}

\usepackage[backend=biber,style=authoryear,citestyle=authoryear, natbib=true]{biblatex}
\addbibresource{papercitations.bib}
\usepackage{setspace}
\onehalfspacing
\addtolength{\skip\footins}{2pc plus 5pt}

\usepackage{geometry}
 \geometry{,
 left=25.4mm,
 top=25.4mm,
 right=25.4mm,
 bottom=25.4mm
 }

 \renewcommand{\TPTnoteSettings}{\footnotesize}
 

\title{Labor Markets and Technological Change: Evidence from Electronic Health Records}


\DeclareLabeldate[online]{%
  \field{date}
  \field{year}
  \field{eventdate}
  \field{origdate}
  \field{urldate}
}
\usepackage{times}


\begin{document}

\linespread{1.2}\title{\vspace{-0.5in} Labor Markets and Technological Change: \\ \Large Evidence from Electronic Health Records\thanks{I thank Ian McCarthy for his continued advice and guidance, along with the additional faculty and graduate students at Emory University for their feedback. Additionally, I would like to thank discussants and conference attendees at the Young Economist Symposium and American Society of Health Economics for their helpful suggestions. I have no conflict of interest or funding sources to declare.}}  
	
	%\date{\today}
	
	\author{\vspace{10mm}Hanna Glenn\footnote{School of Economics, University of Queensland, hanna.glenn@uq.edu.au.}}
    
	
	\maketitle
	%\setlength{\droptitle}{-10pt}
	
	\vspace{-0.2in}
	
	\singlespacing\maketitle

\begin{abstract}\small
    Technology adoption in healthcare has the potential to improve patient outcomes, and in some cases, decrease costs. However, new technology can also impose new burdens on providers. In this paper, I investigate changes in physician labor market behavior following the adoption of electronic health records (EHRs), a major technological shift in healthcare spanning 2009--2015. Treating hospital EHR adoption as an exogenous shock to hospitalists, I use staggered difference-in-differences to estimate the effect of EHRs on hospitalists' decisions to remain in clinical practice, their work setting, and their productivity. I show that hospitalists of retirement age are more likely to leave clinical work after EHR exposure, and those who remain in clinical practice shift work towards office settings. For hospitalists that do not shift towards office work, EHRs lead to a modest increase in patients seen. These findings inform the full costs and benefits of technology adoption in healthcare, which are especially important in a rapidly evolving technological environment.
\end{abstract}

\vspace{2in}
	
	\noindent Keywords: Electronic Health Record, Physician Behavior, Technology Adoption

    \vspace{2mm}
	
	\noindent JEL Codes: I11, J24, O33

	
	\onehalfspacing
	
	\newpage

    \end{document}